\chapterbackmatter{Solutions to Weinberg Exercises}
\ExerciseEditorialNote

\begin{enumerate}
\WeinbergSolution{1}
For bosons, substitute the partition expansion
Eq.~\eqref{eq:4.3.2} into \(F[v]\).  A term containing \(r_k\) connected
components of type \(k\) has the product of their connected matrix
elements and the combinatorial factor \(1/r_k!\) for each set of
identical components.  The factors \(N!M!\) in the definition of \(F\)
cancel the ways of distributing the outgoing and incoming labels among
the components.  Thus
\[
 F[v]
 =\prod_k\sum_{r_k=0}^{\infty}
 \frac{\bigl(F_k^C[v]\bigr)^{r_k}}{r_k!},
\]
where \(F_k^C\) denotes one possible connected contribution, including
all its label integrals.  Multiplying the exponentials and summing over
all component types gives
\[
 \boxed{\qquad F[v]=\exp\!\bigl(F^C[v]\bigr),\qquad
 F^C[v]=\log F[v].\qquad}
\]
The constant term is correct because \(F[0]=1\) and \(F^C[0]=0\).

One can also verify the result by functional differentiation.  Every
derivative of \(\exp(F^C)\) is a sum over set partitions of the
derivatives of \(F^C\); after setting \(v=0\), these derivatives are the
connected matrix elements, and the sum over partitions is exactly
Eq.~\eqref{eq:4.3.2}.

\WeinbergSolution{2}
Let
\[
 \ket{i}=\ket{\mathbf k,-\mathbf k},\qquad
 \ket{f}=\ket{\mathbf k',-\mathbf k'},\qquad
 \varepsilon_k=\sqrt{k^2+M^2}.
\]
Commuting the two annihilation and two creation operators through the
external states gives
\[
 \bra{f}H_{\rm int}\ket{i}
 =4g\,\delta^3(\mathbf p_f-\mathbf p_i).
\]
At second order, the two-particle completeness relation contributes
\(1/2!\).  The reduced transition matrix element is therefore
\[
 \mathcal M(k)
 =4g+8g^2 I(k)+O(g^3),
\qquad
 I(k)=\int d^3\mathbf q\,
 \frac{1}{2\varepsilon_k-2\varepsilon_q+i0}.
\]
The integral is ultraviolet divergent, as expected for an unregulated
three-dimensional contact interaction.  With a momentum cutoff
\(\Lambda\),
\[
 I_\Lambda(k)=4\pi\int_0^\Lambda dq\,
 \frac{q^2}{2\varepsilon_k-2\varepsilon_q+i0},
\qquad
 \operatorname{Im}I_\Lambda(k)=-2\pi^2k\varepsilon_k
\]
for \(\Lambda>k\).  The divergent real part must be absorbed into a
definition of the physical coupling.

Including the identity term, the center-of-mass \(S\)-matrix through
order \(g^2\) is
\[
 \begin{split}
 \OutBra{\mathbf k',-\mathbf k'}\InKet{\mathbf k,-\mathbf k}
 ={}&\braket{\mathbf k',-\mathbf k'}{\mathbf k,-\mathbf k}\\
 &-2\pi i\,\delta(2\varepsilon_{k'}-2\varepsilon_k)
 \delta^3(\mathbf p_f-\mathbf p_i)\,
 \bigl[4g+8g^2I_\Lambda(k)\bigr]+O(g^3).
 \end{split}
\]
Equation \eqref{eq:3.4.30}, with the usual \(1/2!\) for identical final
particles when the full labeled solid angle is used, gives
\[
 \frac{d\sigma}{d\Omega}
 =\frac{(2\pi)^4\varepsilon_k^2}{8}
 \left|4g+8g^2I_\Lambda(k)\right|^2.
\]
In particular, the leading cross section is
\[
 \boxed{\quad
 \frac{d\sigma}{d\Omega}
 =2(2\pi)^4g^2\varepsilon_k^2+O(g^3)
 \quad}
\]
over the full sphere.  Equivalently, one may omit the \(1/2!\) and
integrate over only one hemisphere.

\WeinbergSolution{3}
Put
\[
 \|\lambda\|^2=\int dq\,|\lambda(q)|^2,
\qquad
 A^\dagger(\lambda)=\int dq\,\lambda(q)a_q^\dagger.
\]
The normalized coherent state is
\[
 \boxed{\quad
 \ket{\lambda}
 =e^{-\|\lambda\|^2/2}e^{A^\dagger(\lambda)}\ket{\Omega}
 =e^{-\|\lambda\|^2/2}
 \sum_{N=0}^{\infty}\frac1{N!}
 \int dq_1\cdots dq_N
 \prod_{r=1}^N\lambda(q_r)\ket{q_1\cdots q_N}.
 \quad}
\]
Indeed,
\[
 [a_q,A^\dagger(\lambda)]=\lambda(q)
\]
is a c-number, so
\[
 a_qe^{A^\dagger(\lambda)}
 =e^{A^\dagger(\lambda)}\bigl(a_q+\lambda(q)\bigr).
\]
Since \(a_q\ket{\Omega}=0\), this proves
\(a_q\ket{\lambda}=\lambda(q)\ket{\lambda}\).

Finally,
\[
 \bra{\Omega}e^{A(\lambda)}e^{A^\dagger(\lambda)}
 \ket{\Omega}=e^{\|\lambda\|^2},
\]
so the prefactor makes the norm unity.  The probability of finding
\(N\) particles is
\[
 P_N=e^{-\|\lambda\|^2}\frac{\|\lambda\|^{2N}}{N!},
\]
which also checks the normalization by summing the Poisson series.
\end{enumerate}
